% =====================================================================
% 01-introduction.tex  (IEEE Conference format — no subsections)
% =====================================================================
\section{Introduction}\label{sec:introduction}

\IEEEPARstart{R}{einforcement} Learning (RL) has become the defining
mechanism of the Large Language Model (LLM) post-training pipeline.
After pre-training by next-token prediction and optional Supervised
Fine-Tuning (SFT), it is RL that transforms a fluent yet undirected text
generator into a helpful, harmless, and increasingly \emph{agentic}
system. Reinforcement Learning from Human Feedback
(RLHF)~\cite{christiano2017deeprl,ouyang2022instructgpt} enabled the
first generation of instruction-following assistants; Direct Preference
Optimization (DPO)~\cite{rafailov2023dpo} demonstrated that the same
alignment objective could be achieved without an explicit reward model;
and Group Relative Policy Optimization (GRPO) with verifiable rewards
produced models whose reasoning capabilities are substantially shaped by
RL~\cite{shao2024deepseekmath,guo2025deepseekr1}. Most recently, this
machinery has been extended beyond single responses to optimize
multi-step \emph{trajectories} in which an LLM plans, invokes tools,
maintains memory, and acts in dynamic
environments~\cite{zhang2026landscape}.

This rapid progress, however, has outpaced systematization. Terminology
overlaps across paradigms, algorithms proliferate with variants such as
Proximal Policy Optimization (PPO), DPO, GRPO, and REINFORCE
Leave-One-Out (RLOO), and reward designs range from scalar human
preferences to dense step-level verifiers. Several surveys address parts
of this landscape: some focus on RLHF
alignment~\cite{kaufmann2024survey,casper2023open,lambert2024rlhfbook},
others catalogue broad agent
architectures~\cite{wang2024agentsurvey,xi2025agentic,masterman2024landscape},
and a recent comprehensive taxonomy organizes agentic RL by capabilities
and task domains~\cite{zhang2026landscape}. Yet none provides a
\emph{compact, tutorial-oriented} synthesis that places any given method
along a small set of clearly defined axes. Practitioners entering the
field still lack a concise map connecting preference-based fine-tuning to
the current frontier of agentic RL.

This survey fills that gap with a unified treatment of RL for LLM
training, from its inception in preference-based fine-tuning to
multi-step agentic optimization, organized around a novel three-axis
taxonomy. Our contributions are as follows:

\begin{itemize}
\item A unified Markov Decision Process (MDP) formalization of the
      paradigm shift from single-step \emph{Preference-Based
      Reinforcement Fine-Tuning} (PBRFT) to temporally extended
      \emph{agentic RL}, together with a summary of the algorithmic
      spectrum spanning REINFORCE, PPO, DPO, and GRPO.

\item A \emph{compact three-axis taxonomy} that classifies methods by
      (i)~reward signal, (ii)~optimization algorithm, and
      (iii)~training scope (single response vs.\ multi-step
      trajectory). The three axes are nearly orthogonal, so every method
      occupies a unique, well-defined cell.

\item A structured review of representative works along each axis,
      covering reward design, optimization algorithms, and agentic
      trajectories, consolidated in a master placement table.

\item A structured review of representative works along each axis,
      covering reward design, optimization algorithms, and agentic
      trajectories, consolidated in a master placement table.

\item A synthesis of open challenges, including reward hacking, credit
      assignment, sample efficiency, scalable oversight, and evaluation,
      with concrete research directions.
\end{itemize}

\noindent The remainder of this paper is organized as follows.
Section~\ref{sec:background} formalizes the RL framework for LLMs and
recalls core algorithms. Section~\ref{sec:taxonomy} introduces the
three-axis taxonomy. Sections~\ref{sec:reward}, \ref{sec:optimization},
and~\ref{sec:agentic} examine each axis in turn.
Section~\ref{sec:challenges} discusses open problems, and
Section~\ref{sec:conclusion} concludes.
